% GENERATED FILE. Edit canonical lessons, then run npm run generate:books.
% canonical-source-hash: fnv1a64:d4952a00b3bd6900
% canonical-lessons: JA-C312-rekishi, JA-C312-jidai, JA-C312-kako, JA-C312-izen, JA-C312-touji

\chapter{History}
\label{ch:ja-history}
\hlchaptermodality{312}

\begin{chapteropening}
\textbf{By the end of this chapter:} \emph{I can say what a place was like in the past.}

Complete the last lesson of chapter 312: I can say what a place was like in the past.
\end{chapteropening}

\section[\texorpdfstring{\ja{れきし} (rekishi)}{れきし (rekishi)}]{\ja{れきし} (rekishi) — history}

\label{lesson:JA-C312-rekishi}



\begin{quote}
Before the new one: say the Japanese for to lose for good, then the Japanese for to look back (on).
\end{quote}

\subsection*{You'll want to know: \ja{れきし}}
\textbf{\ja{れきし}} — \emph{rekishi} — \textquotedblleft{}history\textquotedblright{}.

\begin{quote}
\textbf{\ja{日本}\ja{の} \ja{れきしを} \ja{ならっています。}} — \emph{nihon no rekishi o naratte imasu} — \textquotedblleft{}I am learning Japanese history.\textquotedblright{}
\end{quote}

\begin{cousinweb}
\textbf{\ja{れきし}} (\emph{rekishi}) is Sino-Japanese (kango): two Chinese characters read with their on'yomi, readings borrowed from Chinese. They are \emph{reki}, \textquotedblleft{}passage of time\textquotedblright{} (Middle Chinese \emph{lek}) and \emph{shi}, \textquotedblleft{}history\textquotedblright{} (\emph{sriX}) (KANJIDIC2; Middle Chinese in Baxter's transcription).
\end{cousinweb}

\begin{grammarlens}[title={What you've built}]
One more word for telling what happened.
\end{grammarlens}

\subsection*{Your turn}
\begin{itemize}
\raggedright
  \item \emph{Say it:} \emph{rekishi}
  \item \emph{Say it:} \emph{rekishi}, once more
  \item \emph{Say it:} \emph{nihon no rekishi o naratte imasu} — I am learning Japanese history.
  \item \emph{Recall:} say the Japanese for to lose for good, then the Japanese for to look back (on), then say \emph{rekishi} again
\end{itemize}

\begin{tcolorbox}[breakable,colback=teal!4,colframe=teal!35!black,title={Before you move on}]
What does \ja{れきし} mean? (\textquotedblleft{}History\textquotedblright{}.) Say it once more.
\end{tcolorbox}
\section[\texorpdfstring{\ja{じだい} (jidai)}{じだい (jidai)}]{\ja{じだい} (jidai) — an era; a period}

\label{lesson:JA-C312-jidai}



\begin{quote}
Before the new one: say the Japanese for to look back (on), then the Japanese for history.
\end{quote}

\subsection*{You'll want to know: \ja{じだい}}
\textbf{\ja{じだい}} — \emph{jidai} — \textquotedblleft{}an era; a period\textquotedblright{}.

\begin{quote}
\textbf{\ja{さむらいの} \ja{じだいは} \ja{ながかったです。}} — \emph{samurai no jidai wa nagakatta desu} — \textquotedblleft{}The age of the samurai was long.\textquotedblright{}
\end{quote}

\begin{cousinweb}
\textbf{\ja{じだい}} (\emph{jidai}) is Sino-Japanese (kango): two Chinese characters read with their on'yomi, readings borrowed from Chinese. They are \emph{ji}, \textquotedblleft{}time\textquotedblright{} (Middle Chinese \emph{dzyi}) and \emph{dai}, \textquotedblleft{}age\textquotedblright{} (\emph{dojH}) (KANJIDIC2; Middle Chinese in Baxter's transcription).
\end{cousinweb}

\begin{grammarlens}[title={What you've built}]
One more word for telling what happened.
\end{grammarlens}

\subsection*{Your turn}
\begin{itemize}
\raggedright
  \item \emph{Say it:} \emph{jidai}
  \item \emph{Say it:} \emph{jidai}, once more
  \item \emph{Say it:} \emph{samurai no jidai wa nagakatta desu} — The age of the samurai was long.
  \item \emph{Recall:} say the Japanese for to look back (on), then the Japanese for history, then say \emph{jidai} again
\end{itemize}

\begin{tcolorbox}[breakable,colback=teal!4,colframe=teal!35!black,title={Before you move on}]
What does \ja{じだい} mean? (\textquotedblleft{}An era; a period\textquotedblright{}.) Say it once more.
\end{tcolorbox}
\section[\texorpdfstring{\ja{かこ} (kako)}{かこ (kako)}]{\ja{かこ} (kako) — the past}

\label{lesson:JA-C312-kako}



\begin{quote}
Before the new one: say the Japanese for history, then the Japanese for an era.
\end{quote}

\subsection*{You'll want to know: \ja{かこ}}
\textbf{\ja{かこ}} — \emph{kako} — \textquotedblleft{}the past\textquotedblright{}.

\begin{quote}
\textbf{\ja{かこの} \ja{じけんを} \ja{しらべました。}} — \emph{kako no jiken o shirabemashita} — \textquotedblleft{}I looked into incidents from the past.\textquotedblright{}
\end{quote}

\begin{cousinweb}
\textbf{\ja{かこ}} (\emph{kako}) is Sino-Japanese (kango): two Chinese characters read with their on'yomi, readings borrowed from Chinese. They are \emph{ka}, \textquotedblleft{}go beyond\textquotedblright{} (Middle Chinese \emph{kwaH}) and \emph{ko}, \textquotedblleft{}past\textquotedblright{} (\emph{khjoH}) (KANJIDIC2; Middle Chinese in Baxter's transcription).
\end{cousinweb}

\begin{grammarlens}[title={What you've built}]
One more word for telling what happened.
\end{grammarlens}

\subsection*{Your turn}
\begin{itemize}
\raggedright
  \item \emph{Say it:} \emph{kako}
  \item \emph{Say it:} \emph{kako}, once more
  \item \emph{Say it:} \emph{kako no jiken o shirabemashita} — I looked into incidents from the past.
  \item \emph{Recall:} say the Japanese for history, then the Japanese for an era, then say \emph{kako} again
\end{itemize}

\begin{tcolorbox}[breakable,colback=teal!4,colframe=teal!35!black,title={Before you move on}]
What does \ja{かこ} mean? (\textquotedblleft{}The past\textquotedblright{}.) Say it once more.
\end{tcolorbox}
\section[\texorpdfstring{\ja{いぜん} (izen)}{いぜん (izen)}]{\ja{いぜん} (izen) — formerly; in the past}

\label{lesson:JA-C312-izen}



\begin{quote}
Before the new one: say the Japanese for an era, then the Japanese for the past.
\end{quote}

\subsection*{You'll want to know: \ja{いぜん}}
\textbf{\ja{いぜん}} — \emph{izen} — \textquotedblleft{}formerly; in the past\textquotedblright{}.

\begin{quote}
\textbf{\ja{いぜん、ここに} \ja{がっこうが} \ja{ありました。}} — \emph{izen, koko ni gakkō ga arimashita} — \textquotedblleft{}There used to be a school here.\textquotedblright{}
\end{quote}

\begin{cousinweb}
\textbf{\ja{いぜん}} (\emph{izen}) is Sino-Japanese (kango): two Chinese characters read with their on'yomi, readings borrowed from Chinese. They are \emph{i}, \textquotedblleft{}by means of\textquotedblright{} (Middle Chinese \emph{yiX}) and \emph{zen}, \textquotedblleft{}before\textquotedblright{} (\emph{dzen}) (KANJIDIC2; Middle Chinese in Baxter's transcription).
\end{cousinweb}

\begin{grammarlens}[title={What you've built}]
One more word for telling what happened.
\end{grammarlens}

\subsection*{Your turn}
\begin{itemize}
\raggedright
  \item \emph{Say it:} \emph{izen}
  \item \emph{Say it:} \emph{izen}, once more
  \item \emph{Say it:} \emph{izen, koko ni gakkō ga arimashita} — There used to be a school here.
  \item \emph{Recall:} say the Japanese for an era, then the Japanese for the past, then say \emph{izen} again
\end{itemize}

\begin{tcolorbox}[breakable,colback=teal!4,colframe=teal!35!black,title={Before you move on}]
What does \ja{いぜん} mean? (\textquotedblleft{}Formerly; in the past\textquotedblright{}.) Say it once more.
\end{tcolorbox}
\section[\texorpdfstring{\ja{とうじ} (tōji)}{とうじ (tōji)}]{\ja{とうじ} (tōji) — at that time; in those days}

\label{lesson:JA-C312-touji}



\begin{quote}
Before the new one: say the Japanese for the past, then the Japanese for formerly.
\end{quote}

\subsection*{You'll want to know: \ja{とうじ}}
\textbf{\ja{とうじ}} — \emph{tōji} — \textquotedblleft{}at that time; in those days\textquotedblright{}.

\begin{quote}
\textbf{\ja{とうじ、ちちは} \ja{まだ} \ja{がくせいでした。}} — \emph{tōji, chichi wa mada gakusei deshita} — \textquotedblleft{}In those days my father was still a student.\textquotedblright{}
\end{quote}

\begin{cousinweb}
\textbf{\ja{とうじ}} (\emph{tōji}) is Sino-Japanese (kango): two Chinese characters read with their on'yomi, readings borrowed from Chinese. They are \emph{tō}, \textquotedblleft{}right\textquotedblright{} (Middle Chinese \emph{tang}) and \emph{ji}, \textquotedblleft{}time\textquotedblright{} (\emph{dzyi}) (KANJIDIC2; Middle Chinese in Baxter's transcription).
\end{cousinweb}

\begin{grammarlens}[title={What you've built}]
One more word for telling what happened.
\end{grammarlens}

\subsection*{Your turn}
\begin{itemize}
\raggedright
  \item \emph{Say it:} \emph{tōji}
  \item \emph{Say it:} \emph{tōji}, once more
  \item \emph{Say it:} \emph{tōji, chichi wa mada gakusei deshita} — In those days my father was still a student.
  \item \emph{Recall:} say the Japanese for the past, then the Japanese for formerly, then say \emph{tōji} again
\end{itemize}

\begin{tcolorbox}[breakable,colback=teal!4,colframe=teal!35!black,title={Before you move on}]
What does \ja{とうじ} mean? (\textquotedblleft{}At that time; in those days\textquotedblright{}.) Say it once more.
\end{tcolorbox}
